% Paper Chi v1 — Wilson-Action Fisher Zeros for SU(4N)
% Author: Grzegorz Olbryk  <g.olbryk@gmail.com>
% March 2026
% Companion to Paper Rho: the N ≡ 0 mod 4 story

\documentclass[12pt]{article}
\usepackage{amsmath,amssymb,amsthm}
\usepackage[margin=2.5cm]{geometry}
\usepackage{hyperref}
\usepackage{booktabs}

\newtheorem{theorem}{Theorem}
\newtheorem{lemma}[theorem]{Lemma}
\newtheorem{proposition}[theorem]{Proposition}
\newtheorem{corollary}[theorem]{Corollary}
\theoremstyle{definition}
\newtheorem{definition}[theorem]{Definition}
\theoremstyle{remark}
\newtheorem{remark}[theorem]{Remark}

\newcommand{\SU}{\mathrm{SU}}
\newcommand{\Tr}{\mathrm{Tr}}
\newcommand{\dU}{\,\mathrm{d}U}
\newcommand{\avg}[1]{\langle #1 \rangle}
\newcommand{\ord}{\mathrm{ord}}

\title{Wilson-Action Fisher Zeros Beyond the Heat Kernel:\\
Infinitude, Conveyor Belt, and Representation Asymptotics\\
for $\SU(N)$ with $N \equiv 0\!\pmod{4}$}
\author{Grzegorz Olbryk\\[-2pt]
\small\texttt{g.olbryk@gmail.com}}
\date{March 2026}

\begin{document}
\maketitle

\begin{abstract}
For $N \equiv 0\!\pmod{4}$, the dominant saddle-point phase of the
$\SU(N)$ Wilson action vanishes ($\Phi_0 = 0$), and the saddle-point
spacing formula $\avg{\Delta} = \pi/(n|\Phi_0|)$ breaks down.
We prove that the partition function nonetheless has infinitely many
Fisher zeros (Hadamard--Riemann--Lebesgue argument), identify the
\emph{conveyor belt} mechanism by which zeros are produced through
cancellation between progressively higher representations, and
establish that the Wilson-action character integrals $A_p(\kappa)$
decay super-exponentially as $(e\kappa N/(2p))^p$ --- slower than the
heat-kernel prediction $e^{-\kappa p^2/(2N)}$. We identify three
distinct coupling scales ($\kappa_c^{\mathrm{HK}} < \kappa_c^{\mathrm{W}}
\ll \kappa_c^{\mathrm{GWW}}$) with opposite monotonicity in the
representation dominance ratio, showing that heat-kernel predictions
for Fisher zeros are qualitatively unreliable.
\end{abstract}


%====================================================================
\section{Introduction}
%====================================================================

In the companion paper~[Rho], we proved that for all $N$ with
$|\Phi_0| \neq 0$ (odd~$N$ and $N \equiv 2\!\pmod{4}$), the Fisher
zeros of the $n$-plaquette $\SU(N)$ partition function have asymptotic
mean spacing $\avg{\Delta} = \pi/(n|\Phi_0|)$. The present paper treats
the excluded case $N \equiv 0\!\pmod{4}$.

For this parity class, the dominant balanced-split saddle points have
$\Phi_0 = 0$: the Wilson phase vanishes at the saddle. This means
$A_0(\kappa + iy) \sim v_0/y^{m_N}$ with no oscillation, and the
trivial-representation term $A_0^n$ cannot produce zeros on its own.

The natural question is: \emph{does} the partition function have Fisher
zeros at all? The heat-kernel approximation $A_p^{\mathrm{HK}} = d_p
e^{-C_2(p)\kappa}$ predicts that $A_0 = 1$ (constant), and the
sub-dominant oscillating terms die out above a critical~$y$, suggesting
finitely many zeros. This prediction is \emph{wrong}.

We prove:
\begin{enumerate}
  \item \textbf{Infinitude} (Theorem~\ref{thm:infinitude}): $Z_n(s)$ has
  infinitely many zeros for ALL $N \geq 2$.
  \item \textbf{Conveyor belt} (\S\ref{sec:conveyor}): zeros arise from
  cancellation between $A_0^n$ and $d_{p^*} A_{p^*}^n$ where $p^*$
  increases with~$y$.
  \item \textbf{Super-exponential decay} (Theorem~\ref{thm:decay}):
  $A_p(\kappa)/d_p \sim (e\kappa N/(2p))^p$, slower than heat-kernel.
  \item \textbf{Three-scale hierarchy} (\S\ref{sec:hierarchy}):
  heat-kernel and Wilson actions have opposite monotonicity in $A_1/A_0$
  near $\kappa \sim 1$.
\end{enumerate}


%====================================================================
\section{Infinitude Theorem}\label{sec:infinitude}
%====================================================================

\begin{theorem}[Infinitude of Fisher zeros]\label{thm:infinitude}
For all $N \geq 2$ and $n \geq 1$, the Wilson-action partition function
$Z_n(s) = \sum_\lambda d_\lambda A_\lambda(s)^n$ has infinitely many
zeros in~$\mathbb{C}$.
\end{theorem}

\begin{proof}
\textbf{Step 1.} $Z_n(s)$ is entire of exponential type.
Since $|A_\lambda(s)| \leq e^{|\mathrm{Re}\,s| \cdot N}$ and
$\sum_\lambda d_\lambda^2 = N!/(1! \cdots (N-1)!)$ (Plancherel for finite
truncation, convergent for full sum by super-exponential decay of $A_\lambda$),
$Z_n$ is entire with $|Z_n(s)| \leq C \cdot e^{n N |\mathrm{Re}\,s|}$.
Its order is at most~1.

\textbf{Step 2.} $Z_n(\kappa + iy) \to 0$ as $|y| \to \infty$.
Each $A_\lambda(\kappa + iy) = \int_{\SU(N)} \chi_\lambda(U^{-1})
e^{\kappa\,\mathrm{Re}\Tr U} \cdot e^{iy\,\mathrm{Re}\Tr U}\dU$.
The function $g(U) = \chi_\lambda(U^{-1}) e^{\kappa\,\mathrm{Re}\Tr U}$
is in $L^1(\SU(N))$, and $\mathrm{Re}\Tr U$ is a non-constant
continuous function on the compact group. By the Riemann--Lebesgue lemma:
$A_\lambda(\kappa + iy) \to 0$ as $|y| \to \infty$.
Since $Z_n$ is a convergent sum (super-exponential decay gives uniform
convergence), $Z_n(\kappa + iy) \to 0$.

\textbf{Step 3.} By Hadamard's factorisation theorem, an entire function
$f$ of order $\leq 1$ with finitely many zeros is $f(s) = P(s) e^{\alpha s + \beta}$
where $P$ is a polynomial. But $|P(s)e^{\alpha s + \beta}|$ cannot decay
to zero as $|y| \to \infty$ in both the $y > 0$ and $y < 0$ directions
(since $\mathrm{Re}(\alpha(\kappa + iy)) = \alpha\kappa$ is $y$-independent,
$|e^{\alpha s}|$ is constant along vertical lines). Contradiction.
\end{proof}


%====================================================================
\section{The Conveyor Belt Mechanism}\label{sec:conveyor}
%====================================================================

We restrict to $\SU(4)$ ($N \equiv 0 \bmod 4$) as the primary example.
The character expansion gives $Z_n(s) = \sum_{p=0}^\infty d_p A_p(s)^n$,
where the sum is over symmetric-power representations (the non-symmetric
representations contribute to a parallel but smaller sum).

\begin{definition}[Partner representation]
At a zero $s_0 = \kappa + iy_0$ of $Z_n$, the \emph{partner} $p^*$ is
the representation minimising $|A_0(s_0)^n + d_{p^*} A_{p^*}(s_0)^n|$
subject to $d_{p^*} |A_{p^*}(s_0)|^n \geq |A_0(s_0)|^n / 10$.
\end{definition}

\noindent
\textbf{Numerical observation (SU(4), $n = 3$, $\kappa = 1$).}
At the first 14 positive zeros of $Z_3(\kappa + iy)$, the partner
sequence is $p^* = 1, 3, 5, 7, 9, 10, 12, 14, 16, 18, 20, 22, 24, 26$,
increasing approximately linearly with zero index. The phase difference
$|\arg(A_0^n) - \arg(d_{p^*} A_{p^*}^n)|$ is within $0.3$~radians of~$\pi$
at each zero, confirming near-exact anti-phase cancellation.

\begin{remark}
The trivial-representation term $A_0^3$ has total phase accumulation of
only $-0.17$~radians from $y = 1$ to $y = 11$ (approximately $0.03$ turns).
It is essentially a real, decaying function. The zeros are NOT produced by
$A_0$ oscillating past zero; they are produced by a partner $T_{p^*}$
that is in anti-phase with $T_0$.
\end{remark}


%====================================================================
\section{Two-Mechanism Classification}\label{sec:two-mech}
%====================================================================

Fisher zeros of $Z_n(s)$ arise from one of two mechanisms:

\medskip
\begin{center}
\begin{tabular}{lll}
\toprule
& \textbf{Mechanism 1: Saddle-point} & \textbf{Mechanism 2: Conveyor belt} \\
\midrule
Condition & $|\Phi_0| \neq 0$ & $|\Phi_0| = 0$ \\
$N$ class & Odd, $\equiv 2\!\pmod{4}$ & $\equiv 0\!\pmod{4}$ \\
Source & $A_0$ oscillates at frequency $|\Phi_0|$ & $A_0$ real, non-oscillating \\
Cancellation & Within $A_0^n$ itself & Between $T_0$ and $T_{p^*}$ \\
Spacing & Exact: $\Delta y = \pi/(n|\Phi_0|)$ & No closed form \\
Partner & Fixed (saddle pair) & Shifts: $p^* \propto y$ \\
\bottomrule
\end{tabular}
\end{center}


%====================================================================
\section{Large-$p$ Asymptotics}\label{sec:asymptotics}
%====================================================================

\begin{theorem}[Super-exponential decay]\label{thm:decay}
For all $N \geq 2$ and $\kappa > 0$, the Wilson-action character
integrals satisfy
\begin{equation}\label{eq:decay}
  |A_p(\kappa)| \leq C(N)\, p^{N/2 - 1}\,
  \left(\frac{e\kappa N}{2p}\right)^p \qquad \text{for } p > e\kappa N/2.
\end{equation}
Consequently $\sum_{p \geq 0} A_p(\kappa)$ converges, and the generating
function $G(t,\kappa) = \sum_{p \geq 0} A_p(\kappa)\, t^p$ has infinite
radius of convergence.
\end{theorem}

\begin{proof}
By the Weyl integration formula and the Bessel expansion
$e^{\kappa\cos\theta} = \sum_n I_n(\kappa)\, e^{in\theta}$, the
coefficient $A_p$ is a linear combination of products
$\prod_{j=1}^N I_{n_j}(\kappa)$ with $\sum n_j = p + O(N^2)$.
The standard bound \cite{Watson1944}:
$I_n(\kappa) \leq (e\kappa/(2n))^n / \sqrt{2\pi n}$ for $n \geq 1$.
The dominant products have $n_j \approx p/N$, giving
$\prod I_{p/N}(\kappa) \leq (e\kappa N/(2p))^p / p^{N/2}$.
Summing over $O(p^{N-1})$ compositions yields~\eqref{eq:decay}.
Since $(e\kappa N/(2p))^p \to 0$ faster than any exponential,
$|A_p|^{1/p} \to 0$ and $R = \infty$.
\end{proof}

\begin{remark}[Wilson vs.\ heat-kernel]
The heat-kernel predicts $A_p^{\mathrm{HK}} = d_p \exp(-\kappa p(p+N)/(2N))$,
giving $\log(A_p/d_p) \propto -p^2$ (quadratic). The Wilson bound
gives $\log(A_p/d_p) \sim -p \log(2p/(e\kappa N))$, which is $p \log p$ type.
Since $p \log p \ll p^2$ for large~$p$, Wilson decays \emph{slower}:
higher representations contribute more than the heat-kernel predicts.
This makes the conveyor belt possible.
\end{remark}


%====================================================================
\section{Three-Scale Hierarchy}\label{sec:hierarchy}
%====================================================================

For $\SU(N)$ with Wilson action, we identify three coupling scales with
qualitatively different representation structure:

\begin{enumerate}
\item $\kappa_c^{\mathrm{HK}} = 2N\ln N / (N^2 - 1)$:
  the heat-kernel eigenvalue crossing ($A_1^{\mathrm{HK}} = A_0^{\mathrm{HK}}$).
  Decreases with~$N$; $\kappa_c^{\mathrm{HK}} \to 0$.

\item $\kappa_c^{\mathrm{W}}$: the Wilson-action crossing
  ($A_1(\kappa) = A_0(\kappa)$). Increases with~$N$; limit $\approx 2.4$.

\item $\kappa_{\mathrm{half}} \approx 0.85\,N$:
  the Gross--Witten--Wadia eigenvalue concentration scale.
  Linear in~$N$.
\end{enumerate}

\begin{center}
\begin{tabular}{lccc}
\toprule
$N$ & $\kappa_c^{\mathrm{HK}}$ & $\kappa_c^{\mathrm{W}}$ &
$\kappa_{\mathrm{half}}$ \\
\midrule
3 & 0.824 & 1.572 & 2.440 \\
4 & 0.739 & 1.814 & 3.549 \\
5 & 0.671 & 1.941 & 4.622 \\
\bottomrule
\end{tabular}
\end{center}

\noindent
\textbf{Opposite monotonicity.} The heat-kernel ratio
$A_1^{\mathrm{HK}}/A_0^{\mathrm{HK}} = d_1 \exp(-C_2(1)\kappa) = N e^{-(N+1)\kappa/(2N)}$
starts at $d_1 = N$ for $\kappa = 0$ and decreases monotonically.
The Wilson ratio $A_1/A_0$ starts at $0$ (from the character integral)
and \emph{increases} monotonically, crossing~$1$ at $\kappa_c^{\mathrm{W}}$.
These opposite behaviours invalidate heat-kernel predictions for
representation dominance and Fisher zero structure near $\kappa \sim 1$.


%====================================================================
\section{Discussion}
%====================================================================

The $N \equiv 0\!\pmod{4}$ case reveals fundamental differences between
the Wilson and heat-kernel actions. While both are legitimate regularisations
of $\SU(N)$ gauge theory (agreeing in the continuum limit $\kappa \to \infty$),
their analytic structures in the complex-$\kappa$ plane are qualitatively
different. The heat-kernel, with its exact character expansion
$A_p^{\mathrm{HK}} = d_p \exp(-C_2 \kappa)$, has no interaction between
representations. The Wilson action, through the generating function
$G(t,s) = \int e^{s\,\mathrm{Re}\Tr U} / \det(I - tU)\dU$, couples all
representations and produces the conveyor belt mechanism.

\medskip\noindent
\textbf{Open problems.}
\begin{enumerate}
  \item Analytical formula for the conveyor belt spacing.
  \item Universality of the phase progression rate $\beta \approx 1.7\pi$/rep.
  \item The thermodynamic limit for $N \equiv 0\!\pmod{4}$ (zero density
  is $0$ by the saddle-point argument, but the true density from the
  conveyor belt is nonzero).
\end{enumerate}


%====================================================================
\begin{thebibliography}{99}

\bibitem{Fisher1965}
M.~E.~Fisher, \emph{The Nature of Critical Points},
Lectures in Theoretical Physics, vol.~7C (1965).

\bibitem{GW80}
D.~J.~Gross and E.~Witten,
\emph{Possible third-order phase transition in the large-$N$ lattice
gauge theory}, Phys.\ Rev.\ D \textbf{21} (1980) 446.

\bibitem{Pi}
G.~Olbryk, \emph{Asymptotic spacing of Fisher zeros in a two-plaquette
$\SU(N)$ lattice gauge model} (Paper~Pi, v27), preprint (2026).

\bibitem{Rho}
G.~Olbryk, \emph{Unified asymptotic spacing of Fisher zeros in
multi-plaquette $\SU(N)$ lattice gauge models} (Paper~Rho, v1),
preprint (2026).

\bibitem{Watson1944}
G.~N.~Watson, \emph{A Treatise on the Theory of Bessel Functions},
Cambridge University Press (1944).

\bibitem{BtD85}
T.~Br\"ocker and T.~tom~Dieck, \emph{Representations of Compact Lie
Groups}, Springer GTM~98 (1985).

\bibitem{Hormander85}
L.~H\"ormander, \emph{The Analysis of Linear Partial Differential
Operators~I}, Springer (1985).

\bibitem{Wadia80}
S.~Wadia, \emph{$N = \infty$ phase transition in a class of exactly
soluble model lattice gauge theories}, Phys.\ Lett.\ B \textbf{93} (1980) 403.

\end{thebibliography}

\end{document}
